\documentclass[11pt]{article}

%%% PACKAGES

\usepackage{fullpage}  % 1 page margins
\usepackage{latexsym}  % extra symbols, e.g. \leadsto
\usepackage{verbatim}  % verbatim mode
\usepackage{proof}     % proof mode
\usepackage{amsthm,amssymb,amsmath}  % various math symbols
\usepackage{color}     % color control
\usepackage{etoolbox}  % misc utilities
\usepackage{graphics}  % images
\usepackage{mathpartir}% inference rules
\usepackage{hyperref}  % hyperlinks
\usepackage{titlesec}  % title/section controls
\usepackage{minted}    % code blocks
\usepackage[hang,flushmargin]{footmisc} % No indent footnotes
\usepackage{parskip}   % no parindent
\usepackage{tikz}      % drawing figures
\usepackage[T1]{fontenc}
%%% COMMANDS

% Typography and symbols
\newcommand{\msf}[1]{\mathsf{#1}}
\newcommand{\ctx}{\Gamma}
\newcommand{\qamp}{&\quad}
\newcommand{\qqamp}{&&\quad}
\newcommand{\Coloneqq}{::=}
\newcommand{\proves}{\vdash}
\newcommand{\str}[1]{{#1}^{*}}
\newcommand{\eps}{\varepsilon}
\newcommand{\brc}[1]{\{{#1}\}}
\newcommand{\binopm}[2]{{#1}~\bar{\oplus}~{#2}}
\newcommand{\aequiv}{\equiv_\alpha}

% Untyped lambda calculus
\newcommand{\fun}[2]{\lambda ~ {#1} ~ . ~ {#2}}
\newcommand{\app}[2]{{#1} ~ {#2}}
\newcommand{\fix}[1]{\msf{fix}~{#1}}

% Typed lambda calculus - expressions
\newcommand{\funt}[3]{\lambda ~ ({#1} : {#2}) ~ . ~ {#3}}
\newcommand{\ift}[3]{\msf{if} ~ {#1} ~ \msf{then} ~ {#2} ~ \msf{else} ~ {#3}}
\newcommand{\rec}[5]{\msf{rec}(#1; ~ #2.#3.#4)(#5)}
\newcommand{\lett}[4]{\msf{let} ~ \hasType{#1}{#2} = {#3} ~ \msf{in} ~ {#4}}
\newcommand{\falset}{\msf{false}}
\newcommand{\truet}{\msf{true}}
\newcommand{\case}[5]{\msf{case} ~ {#1} ~ \{ L({#2}) \to {#3} \mid R({#4}) \to {#5} \}}
\newcommand{\pair}[2]{({#1},{#2})}
\newcommand{\proj}[2]{{#1} . {#2}}
\newcommand{\inj}[3]{\msf{inj} ~ {#1} = {#2} ~ \msf{as} ~ {#3}}
\newcommand{\letv}[3]{\msf{let} ~ {#1} = {#2} ~ \msf{in} ~ {#3}}
\newcommand{\fold}[2]{\msf{fold}~{#1}~\msf{as}~{#2}}
\newcommand{\unfold}[1]{\msf{unfold}~{#1}}
\newcommand{\poly}[2]{\Lambda~{#1}~.~{#2}}
\newcommand{\polyapp}[2]{{#1}~[{#2}]}

% Typed lambda calculus - types
\newcommand{\tnum}{\msf{number}}
\newcommand{\tstr}{\msf{string}}
\newcommand{\tint}{\msf{int}}
\newcommand{\tbool}{\msf{bool}}
\newcommand{\tfun}[2]{{#1} \rightarrow {#2}}
\newcommand{\tprod}[2]{{#1} \times {#2}}
\newcommand{\tsum}[2]{{#1} + {#2}}
\newcommand{\trec}[2]{\mu~{#1}~.~{#2}}
\newcommand{\tvoid}{\msf{void}}
\newcommand{\tunit}{\msf{unit}}
\newcommand{\tpoly}[2]{\forall~{#1}~.~{#2}}

% WebAssembly
\newcommand{\wconst}[1]{\msf{i32.const}~{#1}}
\newcommand{\wbinop}[1]{\msf{i32}.{#1}}
\newcommand{\wgetlocal}[1]{\msf{get\_local}~{#1}}
\newcommand{\wsetlocal}[1]{\msf{set\_local}~{#1}}
\newcommand{\wgetglobal}[1]{\msf{get\_global}~{#1}}
\newcommand{\wsetglobal}[1]{\msf{set\_global}~{#1}}
\newcommand{\wload}{\msf{i32.load}}
\newcommand{\wstore}{\msf{i32.store}}
\newcommand{\wsize}{\msf{memory.size}}
\newcommand{\wgrow}{\msf{memory.grow}}
\newcommand{\wunreachable}{\msf{unreachable}}
\newcommand{\wblock}[1]{\msf{block}~{#1}}
\newcommand{\wloop}[1]{\msf{loop}~{#1}}
\newcommand{\wbr}[1]{\msf{br}~{#1}}
\newcommand{\wbrif}[1]{\msf{br\_if}~{#1}}
\newcommand{\wreturn}{\msf{return}}
\newcommand{\wcall}[1]{\msf{call}~{#1}}
\newcommand{\wlabel}[2]{\msf{label}~\{#1\}~{#2}}
\newcommand{\wframe}[2]{\msf{frame}~({#1}, {#2})}
\newcommand{\wtrapping}{\msf{trapping}}
\newcommand{\wbreaking}[1]{\msf{breaking}~{#1}}
\newcommand{\wreturning}[1]{\msf{returning}~{#1}}
\newcommand{\wconfig}[5]{\{\msf{module}{:}~{#1};~\msf{mem}{:}~{#2};~\msf{locals}{:}~{#3};~\msf{stack}{:}~{#4};~\msf{instrs}{:}~{#5}\}}
\newcommand{\wfunc}[4]{\{\msf{params}{:}~{#1};~\msf{locals}{:}~{#2};~\msf{return}~{#3};~\msf{body}{:}~{#4}\}}
\newcommand{\wmodule}[1]{\{\msf{funcs}{:}~{#1}\}}
\newcommand{\wcg}{\msf{globals}}
\newcommand{\wcf}{\msf{funcs}}
\newcommand{\wci}{\msf{instrs}}
\newcommand{\wcs}{\msf{stack}}
\newcommand{\wcl}{\msf{locals}}
\newcommand{\wcm}{\msf{mem}}
\newcommand{\wcmod}{\msf{module}}
\newcommand{\wsteps}[2]{\steps{\brc{#1}}{\brc{#2}}}
\newcommand{\with}{\underline{\msf{with}}}
\newcommand{\wvalid}[2]{{#1} \vdash {#2}~\msf{valid}}
\newcommand{\wif}[2]{\msf{if}~{#1}~{\msf{else}}~{#2}}
\newcommand{\wfor}[4]{\msf{for}~(\msf{init}~{#1})~(\msf{cond}~{#2})~(\msf{post}~{#3})~{#4}}
% assign4.3 custom
\newcommand{\wtry}[2]{\msf{try}~{#1}~\msf{catch}~{#2}}
\newcommand{\wraise}{\msf{raise}}
\newcommand{\wraising}[1]{\msf{raising}~{#1}}

% Inference rules
%\newcommand{\inferrule}[3][]{\cfrac{#2}{#3}\;{#1}}
\newcommand{\ir}[3]{\inferrule*[right=\text{(#1)}]{#2}{#3}}
\newcommand{\s}{\hspace{1em}}
\newcommand{\nl}{\\[2em]}
\newcommand{\steps}[2]{#1 \boldsymbol{\mapsto} #2}
\newcommand{\evals}[2]{#1 \boldsymbol{\overset{*}{\mapsto}} #2}
\newcommand{\subst}[3]{[#1 \rightarrow #2] ~ #3}
\newcommand{\dynJ}[2]{#1 \proves #2}
\newcommand{\dynJC}[1]{\dynJ{\ctx}{#1}}
\newcommand{\typeJ}[3]{#1 \proves \hasType{#2}{#3}}
\newcommand{\typeJC}[2]{\typeJ{\ctx}{#1}{#2}}
\newcommand{\hasType}[2]{#1 : #2}
\newcommand{\val}[1]{#1~\msf{val}}
\newcommand{\num}[1]{\msf{Int}(#1)}
\newcommand{\err}[1]{#1~\msf{err}}
\newcommand{\trans}[2]{#1 \leadsto #2}
\newcommand{\size}[1]{|#1|}

\newcommand{\hwtitle}[2]{\begin{center}{\Large #1} \\[0.5em] {\large #2}\end{center}\vspace{1em}}
\newcommand{\toc}{{\hypersetup{hidelinks}\tableofcontents}}
\newcommand{\problem}[1]{\section*{#1}}

%%% CONFIG

% Spacing around title/sections
\titlelabel{\thetitle.\quad}
\titlespacing*{\section}{0pt}{10pt}{0pt}
\titlespacing*{\subsection}{0pt}{10pt}{0pt}


% Show color on hyperlinks
\hypersetup{colorlinks=true}

% Stylize code blocks
\usemintedstyle{xcode}
% TODO(wcrichto): remove line of spacing beneath code blocks
\newcommand{\nm}[2]{
  \newminted{#1}{#2}
  \newmint{#1}{#2}
  \newmintinline{#1}{#2}}
\nm{lua}{}
\nm{ocaml}{}
\nm{rust}{}
\nm{prolog}{}

\newcommand{\ml}[1]{\ocamlinline|#1|}


\begin{document}

\hwtitle
  {Assignment 2}
  {Will Crichton (wcrichto)} %% REPLACE THIS WITH YOUR NAME/ID

\problem{Problem 2}

\textbf{Part 1:}

\begin{mathpar}
\text{Step 1:}\qquad
\ir{D-App-Body}
  {\ir{D-App-Lam}
    {\ir{D-App-Done}
      {\ir{D-Lam}{ \ }{\val{\fun{\_}{x}}}}
      {\dynJ{\{x \rightarrow D\}}{\steps
        {\app{(\fun{x}{\fun{\_}{x}})}{L}}
        {\fun{\_}{x}}}}}
    {\dynJ{\{x \rightarrow D\}}{\steps
      {\app{\app{(\fun{x}{\fun{\_}{x}})}{L}}{*}}
      {\app{(\fun{\_}{x})}{*}}}}}
  {\dynJ{\varnothing}{\steps
    {\app{(\fun{x}{\app{\app{(\fun{x}{\fun{\_}{x}})}{L}}{*}})}{D}}
    {\app{(\fun{x}{\app{(\fun{\_}{x})}{*}})}{D}}}}
\end{mathpar}

Fill in the remaining steps here.

\begin{mathpar}
\text{Step 2:}\qquad
\ir{D-App-Body}
  {\ir{D-App-Body}
    {\ir{D-Var}
      {x \rightarrow D \in \{x \rightarrow D, \_ \rightarrow *\}}
      {\dynJ{\{x \rightarrow D, \_ \rightarrow *\}}
        {\steps{x}{D}}}}
    {\dynJ{\{x \rightarrow D\}}{\steps
      {\app{(\fun{\_}{x})}{*}}
      {\app{(\fun{\_}{D})}{*}}}}}
  {\dynJ{\varnothing}{\steps
    {\app{(\fun{x}{\app{(\fun{\_}{x})}{*}})}{D}}
    {\app{(\fun{x}{\app{(\fun{\_}{D})}{*}})}{D}}}}
\end{mathpar}

\begin{mathpar}
\text{Step 3:}\qquad
\ir{D-App-Body}
  {\ir{D-App-Done}
    {\val{D}}
    {\dynJ{\{x \rightarrow D\}}{\steps
      {\app{(\fun{\_}{D})}{*}}
      {D}}}}
  {\dynJ{\varnothing}{\steps
    {\app{(\fun{x}{\app{(\fun{\_}{D})}{*}})}{D}}
    {\app{(\fun{x}{D})}{D}}}}
\end{mathpar}

\begin{mathpar}
\text{Step 4:}\qquad
\ir{D-App-Done}
  {\val{D}}
  {\dynJ{\varnothing}{\steps
    {\app{(\fun{x}{D})}{D}}
    {D}}}
\end{mathpar}

\textbf{Part 2:}

Here's examples of how to LaTeX the existing dynamically-scoped lambda calculus rules. Delete them before you submit.

\begin{mathpar}
\ir{D-Lam}{ \ }{\val{\fun{x}{e}}}

\ir{D-Var}
  {x \rightarrow e \in \ctx}
  {\dynJC{\steps{x}{e}}}

\ir{D-App-Lam}
  {\dynJC{\steps{e_\msf{lam}}{e_\msf{lam}'}}}
  {\dynJC{\steps{\app{e_\msf{lam}}{e_\msf{arg}}}{\app{e_\msf{lam}'}{e_\msf{arg}}}}}

\ir{D-App-Body}
  {\dynJ{\ctx, x \rightarrow e_\msf{arg}}{\steps{e_\msf{body}}{e_\msf{body}'}}}
  {\dynJC{\steps{\app{(\fun{x}{e_\msf{body}})}{e_\msf{arg}}}{\app{(\fun{x}{e_\msf{body}'})}{e_\msf{arg}}}}} \s

\ir{D-App-Done}
  {\val{e_\msf{body}}}
  {\dynJC{\steps{\app{(\fun{x}{e_\msf{body}})}{e_\msf{arg}}}{e_\msf{body}}}}
\end{mathpar}

\begin{mathpar}
\ir{D-Let-Body}
  {\dynJ{\ctx, x \rightarrow e_\msf{var}}
    {\steps{e_\msf{body}}{e_\msf{body}'}}}
  {\dynJC{\steps
    {\letv{x}{e_\msf{var}}{e_\msf{body}}}
    {\letv{x}{e_\msf{var}}{e_\msf{body}'}}}}

\ir{D-Let-Done}
  {\val{e_\msf{body}}}
  {\dynJC{\steps
    {\letv{x}{e_\msf{var}}{e_\msf{body}}}
    {e_\msf{body}}}}
\end{mathpar}

\newpage

\problem{Problem 3}

\textbf{The \texttt{let} extension is not type safe.}

The typing rule \textsc{T-Let} does not require
$e_\msf{var}$ to have the declared type $\tau_\msf{var}$.
Consider the expression
\[
e =
\lett{x}{\tnum}{\truet}{x}.
\]

This expression has type $\tnum$:
\begin{mathpar}
\ir{T-Let}
  {\ir{T-Var}
    {\hasType{x}{\tnum} \in
      \{\hasType{x}{\tnum}\}}
    {\typeJ{\{\hasType{x}{\tnum}\}}{x}{\tnum}}}
  {\typeJ{\varnothing}
    {\lett{x}{\tnum}{\truet}{x}}
    {\tnum}}
\end{mathpar}

Notice that \textsc{T-Let} never checks the type of $\truet$.
By \textsc{D-Let}, the expression takes one step:
\begin{mathpar}
\ir{D-Let}
  {\,}
  {\steps
    {\lett{x}{\tnum}{\truet}{x}}
    {\subst{x}{\truet}{x}}}
\end{mathpar}

Since $\subst{x}{\truet}{x} = \truet$, we have
\[
\steps
  {\lett{x}{\tnum}{\truet}{x}}
  {\truet}.
\]

However,
\begin{mathpar}
\ir{T-True}
  {\,}
  {\typeJ{\varnothing}{\truet}{\tbool}}
\end{mathpar}
and therefore
\[
\varnothing \nvdash \hasType{\truet}{\tnum}.
\]

Thus, we have an expression $e$ such that
\[
\typeJ{\varnothing}{e}{\tnum},
\qquad
\steps{e}{\truet},
\qquad
\varnothing \nvdash \hasType{\truet}{\tnum}.
\]
Therefore, the \texttt{let} extension violates preservation and is
not type safe.


\bigskip
\textbf{The \texttt{rec} extension is type safe.}

\medskip
\textbf{Progress.}

Suppose
\[
\typeJ{\varnothing}
  {\rec{e_\msf{base}}
    {x_\msf{num}}
    {x_\msf{acc}}
    {e_\msf{acc}}
    {e_\msf{arg}}}
  {\tau}.
\]

By inversion on \textsc{T-Rec}, we have
\[
\typeJ{\varnothing}{e_\msf{arg}}{\tnum},
\]
\[
\typeJ{\varnothing}{e_\msf{base}}{\tau},
\]
and
\[
\typeJ{
  \{\hasType{x_\msf{num}}{\tnum},
    \hasType{x_\msf{acc}}{\tau}\}}
  {e_\msf{acc}}
  {\tau}.
\]

By the induction hypothesis applied to $e_\msf{arg}$, either
$e_\msf{arg}$ is a value or it can take a step.

\begin{enumerate}
  \item Suppose there exists $e_\msf{arg}'$ such that
  \[
  \steps{e_\msf{arg}}{e_\msf{arg}'}.
  \]
  Then, by \textsc{D-Rec-Step},
  \[
  \steps
    {\rec{e_\msf{base}}
      {x_\msf{num}}
      {x_\msf{acc}}
      {e_\msf{acc}}
      {e_\msf{arg}}}
    {\rec{e_\msf{base}}
      {x_\msf{num}}
      {x_\msf{acc}}
      {e_\msf{acc}}
      {e_\msf{arg}'}}.
  \]

  \item Suppose $\val{e_\msf{arg}}$. Since
  \[
  \typeJ{\varnothing}{e_\msf{arg}}{\tnum},
  \]
  the canonical forms lemma implies that
  $e_\msf{arg}=n$ for some natural number $n$.

  There are two cases:
  \begin{enumerate}
    \item If $n=0$, then by \textsc{D-Rec-Base},
    \[
    \steps
      {\rec{e_\msf{base}}
        {x_\msf{num}}
        {x_\msf{acc}}
        {e_\msf{acc}}
        {0}}
      {e_\msf{base}}.
    \]

    \item If $n>0$, then by \textsc{D-Rec-Dec},
    \[
    \steps
      {\rec{e_\msf{base}}
        {x_\msf{num}}
        {x_\msf{acc}}
        {e_\msf{acc}}
        {n}}
      {\subst{x_\msf{num}}{n}{
        \subst{x_\msf{acc}}{
          \rec{e_\msf{base}}
            {x_\msf{num}}
            {x_\msf{acc}}
            {e_\msf{acc}}
            {n-1}}
          {e_\msf{acc}}}}.
    \]
  \end{enumerate}
\end{enumerate}

Therefore, every closed, well-typed \texttt{rec} expression can take
a step. Hence progress holds.


\bigskip
\textbf{Preservation.}

Suppose
\[
\typeJC{
  \rec{e_\msf{base}}
    {x_\msf{num}}
    {x_\msf{acc}}
    {e_\msf{acc}}
    {e_\msf{arg}}
}{\tau}
\]
and
\[
\steps
  {\rec{e_\msf{base}}
    {x_\msf{num}}
    {x_\msf{acc}}
    {e_\msf{acc}}
    {e_\msf{arg}}}
  {e'}.
\]

By inversion on \textsc{T-Rec}, we have
\[
\typeJC{e_\msf{arg}}{\tnum},
\tag{1}
\]
\[
\typeJC{e_\msf{base}}{\tau},
\tag{2}
\]
and
\[
\typeJ{
  \ctx,
  \hasType{x_\msf{num}}{\tnum},
  \hasType{x_\msf{acc}}{\tau}}
  {e_\msf{acc}}
  {\tau}.
\tag{3}
\]

We proceed by cases on the rule used in the reduction.

\begin{enumerate}
  \item \textbf{Case \textsc{D-Rec-Step}.}

  In this case,
  \[
  \steps{e_\msf{arg}}{e_\msf{arg}'}
  \]
  and
  \[
  e' =
  \rec{e_\msf{base}}
    {x_\msf{num}}
    {x_\msf{acc}}
    {e_\msf{acc}}
    {e_\msf{arg}'}.
  \]

  By preservation for $e_\msf{arg}$ and (1),
  \[
  \typeJC{e_\msf{arg}'}{\tnum}.
  \]

  Combining this result with (2) and (3), \textsc{T-Rec} gives
  \[
  \typeJC{
    \rec{e_\msf{base}}
      {x_\msf{num}}
      {x_\msf{acc}}
      {e_\msf{acc}}
      {e_\msf{arg}'}
  }{\tau}.
  \]

  Therefore,
  \[
  \typeJC{e'}{\tau}.
  \]

  \item \textbf{Case \textsc{D-Rec-Base}.}

  In this case, $e_\msf{arg}=0$ and
  \[
  e'=e_\msf{base}.
  \]

  By (2),
  \[
  \typeJC{e_\msf{base}}{\tau}.
  \]

  Therefore,
  \[
  \typeJC{e'}{\tau}.
  \]

  \item \textbf{Case \textsc{D-Rec-Dec}.}

  In this case, $e_\msf{arg}=n$, where $n>0$, and
  \[
  e' =
  [x_\msf{num} \rightarrow n,\,
   x_\msf{acc} \rightarrow
     \rec{e_\msf{base}}
       {x_\msf{num}}
       {x_\msf{acc}}
       {e_\msf{acc}}
       {n-1}]
  \ e_\msf{acc}.
  \]

  Since $n$ and $n-1$ are natural numbers,
  \[
  \typeJC{n}{\tnum}
  \qquad\text{and}\qquad
  \typeJC{n-1}{\tnum}.
  \tag{4}
  \]

  Applying \textsc{T-Rec} to (2), (3), and
  $\typeJC{n-1}{\tnum}$ gives
  \[
  \typeJC{
    \rec{e_\msf{base}}
      {x_\msf{num}}
      {x_\msf{acc}}
      {e_\msf{acc}}
      {n-1}
  }{\tau}.
  \tag{5}
  \]

  From (3), (4), and (5), the simultaneous substitution lemma gives
  \[
  \typeJC{
    [x_\msf{num} \rightarrow n,\,
     x_\msf{acc} \rightarrow
       \rec{e_\msf{base}}
         {x_\msf{num}}
         {x_\msf{acc}}
         {e_\msf{acc}}
         {n-1}]
    \ e_\msf{acc}
  }{\tau}.
  \]

  Therefore,
  \[
  \typeJC{e'}{\tau}.
  \]
\end{enumerate}

All possible reduction rules preserve the type $\tau$. Hence
preservation holds for the \texttt{rec} extension.

Since the \texttt{rec} extension satisfies both progress and
preservation, it is type safe.

\end{document}
